% ============================================================
%  CAPÍTULO 12 – Pruebas y Validación
% ============================================================
\chapter{Pruebas y Validación}

\section{Entorno de pruebas}

Las pruebas de FitLabs se realizaron en el siguiente entorno:

\begin{itemize}
  \item \textbf{Dispositivo físico Android}: Smartphone con Android 12 (API Level 31), utilizado para pruebas de rendimiento y usabilidad real.
  \item \textbf{Emulador Android}: Android Studio Emulator con imagen de Pixel 6 (Android 13) para pruebas de regresión rápida durante el desarrollo.
  \item \textbf{Backend}: Proyecto Supabase en plan Free, base de datos poblada con datos de prueba (usuarios, rutinas, ejercicios, mensajes).
  \item \textbf{Red}: Pruebas realizadas tanto en WiFi (latencia baja) como en datos móviles 4G (latencia variable) para evaluar la resiliencia de las llamadas HTTP.
\end{itemize}

\section{Plan de pruebas}

\subsection{Módulo de Autenticación}

\begin{longtable}{@{}p{1cm}p{5cm}p{5cm}p{2.5cm}@{}}
\toprule
\textbf{ID} & \textbf{Caso de prueba} & \textbf{Resultado esperado} & \textbf{Resultado} \\
\midrule
\endhead
PT-01 & Registro con datos válidos & Usuario creado en Supabase Auth y perfil en tabla \texttt{perfiles} & \textcolor{green!60!black}{\textbf{OK}} \\
PT-02 & Registro con email ya existente & Mensaje de error informativo mostrado al usuario & \textcolor{green!60!black}{\textbf{OK}} \\
PT-03 & Login con credenciales correctas & Redirección a \texttt{/resumen} & \textcolor{green!60!black}{\textbf{OK}} \\
PT-04 & Login con contraseña incorrecta & Mensaje de error sin revelar si el email existe & \textcolor{green!60!black}{\textbf{OK}} \\
PT-05 & Cierre de sesión & Token invalidado, redirección al login & \textcolor{green!60!black}{\textbf{OK}} \\
\bottomrule
\end{longtable}

\subsection{Módulo de Clientes}

\begin{longtable}{@{}p{1cm}p{5cm}p{5cm}p{2.5cm}@{}}
\toprule
\textbf{ID} & \textbf{Caso de prueba} & \textbf{Resultado esperado} & \textbf{Resultado} \\
\midrule
\endhead
PT-06 & Listar clientes del entrenador & Solo se muestran los clientes asociados a ese entrenador (RLS) & \textcolor{green!60!black}{\textbf{OK}} \\
PT-07 & Buscar cliente por nombre & Lista filtrada en tiempo real & \textcolor{green!60!black}{\textbf{OK}} \\
PT-08 & Acceder al detalle de un cliente & Pantalla con datos del cliente correcto & \textcolor{green!60!black}{\textbf{OK}} \\
PT-09 & Ver sesiones pasadas del cliente & Lista de sesiones ordenada por fecha & \textcolor{green!60!black}{\textbf{OK}} \\
\bottomrule
\end{longtable}

\subsection{Módulo de Rutinas}

\begin{longtable}{@{}p{1cm}p{5cm}p{5cm}p{2.5cm}@{}}
\toprule
\textbf{ID} & \textbf{Caso de prueba} & \textbf{Resultado esperado} & \textbf{Resultado} \\
\midrule
\endhead
PT-10 & Buscar ejercicio en la API & Lista de ejercicios con nombre, grupo muscular e imagen & \textcolor{green!60!black}{\textbf{OK}} \\
PT-11 & Añadir ejercicio a la rutina & Ejercicio aparece en la lista de la rutina con parámetros & \textcolor{green!60!black}{\textbf{OK}} \\
PT-12 & Guardar rutina en Supabase & Rutina guardada y ejercicios vinculados en la BD & \textcolor{green!60!black}{\textbf{OK}} \\
PT-13 & Crear rutina sin título & Validación impide el envío, muestra aviso & \textcolor{green!60!black}{\textbf{OK}} \\
\bottomrule
\end{longtable}

\subsection{Módulo de Mensajes}

\begin{longtable}{@{}p{1cm}p{5cm}p{5cm}p{2.5cm}@{}}
\toprule
\textbf{ID} & \textbf{Caso de prueba} & \textbf{Resultado esperado} & \textbf{Resultado} \\
\midrule
\endhead
PT-14 & Listar conversaciones & Solo chats donde participa el usuario actual & \textcolor{green!60!black}{\textbf{OK}} \\
PT-15 & Enviar mensaje de texto & Mensaje aparece en la pantalla del remitente & \textcolor{green!60!black}{\textbf{OK}} \\
PT-16 & Recibir mensaje en tiempo real & El receptor ve el mensaje sin refrescar la pantalla & \textcolor{green!60!black}{\textbf{OK}} \\
PT-17 & Marcar mensajes como leídos & El contador de no leídos se actualiza correctamente & \textcolor{green!60!black}{\textbf{OK}} \\
\bottomrule
\end{longtable}

\section{Resultados y conclusiones de las pruebas}

Todos los casos de prueba definidos en el plan resultaron satisfactorios. Los puntos más críticos, como el aislamiento de datos mediante RLS y la actualización en tiempo real del módulo de mensajes, funcionaron correctamente tanto en emulador como en dispositivo físico con conectividad 4G.

Durante las pruebas se detectaron y corrigieron los siguientes incidentes menores:

\begin{itemize}
  \item \textbf{Incidencia 1}: En el módulo de Mis Clientes, el buscador no filtraba correctamente nombres con tildes. Solución: normalización del texto de búsqueda antes de la comparación.
  \item \textbf{Incidencia 2}: La pantalla de Crear Rutina no limpiaba correctamente el estado al navegar hacia atrás y volver a abrirla. Solución: uso de \texttt{initState} para reiniciar la lista de ejercicios.
\end{itemize}

\section{Pruebas no funcionales}

Además de la validación funcional, se realizaron pruebas no funcionales centradas en rendimiento percibido, robustez ante errores de red, seguridad y usabilidad. Estas pruebas permiten valorar si la aplicación es apta para uso diario por entrenadores con distintos niveles de experiencia técnica.

\subsection{Rendimiento percibido en pantallas críticas}

Se midió el tiempo desde la acción del usuario hasta la visualización útil de datos en las pantallas más relevantes.

\begin{longtable}{@{}p{4cm}p{3cm}p{3cm}p{4cm}@{}}
\toprule
\textbf{Pantalla} & \textbf{WiFi (media)} & \textbf{4G (media)} & \textbf{Observaciones} \\
\midrule
\endhead
Login / Registro & 0.8 s & 1.2 s & Validación local inmediata; latencia asociada a Supabase Auth \\
Resumen del Día & 1.1 s & 1.8 s & Carga inicial de contadores y próximas sesiones \\
Mis Clientes & 1.3 s & 2.0 s & Filtrado local fluido tras primer fetch \\
Mensajes & 1.0 s & 1.6 s & Realtime estable, sin bloqueos de UI \\
Crear Rutina & 1.4 s & 2.1 s & Pantalla más exigente por listado dinámico de ejercicios \\
\bottomrule
\caption{Tiempos de respuesta observados en pruebas manuales de rendimiento.}
\end{longtable}

En todos los casos se mantuvo una experiencia utilizable dentro de los objetivos del proyecto para entorno móvil, especialmente tras aplicar optimizaciones simples de reconstrucción de widgets y cargas diferidas.

\subsection{Pruebas de resiliencia y conectividad}

\begin{longtable}{@{}p{1cm}p{5cm}p{5cm}p{2.5cm}@{}}
\toprule
\textbf{ID} & \textbf{Escenario} & \textbf{Resultado esperado} & \textbf{Resultado} \\
\midrule
\endhead
PNF-01 & Pérdida de red durante login & Mostrar error y permitir reintento sin cierre de app & \textcolor{green!60!black}{\textbf{OK}} \\
PNF-02 & Timeout en carga de clientes & Estado de carga controlado y mensaje informativo & \textcolor{green!60!black}{\textbf{OK}} \\
PNF-03 & Reconexión tras corte breve en mensajes & Reanudación automática de suscripción Realtime & \textcolor{green!60!black}{\textbf{OK}} \\
PNF-04 & Datos vacíos en resumen & Mostrar estado vacío sin errores visuales & \textcolor{green!60!black}{\textbf{OK}} \\
PNF-05 & Doble envío rápido de mensaje & Evitar duplicados visibles por control en UI & \textcolor{green!60!black}{\textbf{OK}} \\
\bottomrule
\caption{Pruebas de resiliencia ante fallos de red y estados límite.}
\end{longtable}

\subsection{Validación de seguridad de acceso a datos}

Se verificó el comportamiento esperado de las políticas RLS de Supabase en operaciones de lectura y escritura para evitar exposición de información entre entrenadores.

\begin{itemize}
  \item Intento de lectura de clientes no asociados: denegado por política.
  \item Intento de acceso a chat ajeno: denegado por política.
  \item Consulta de rutinas no vinculadas al usuario autenticado: sin resultados.
  \item Escritura de mensajes en chat no participante: bloqueada por condición de pertenencia.
\end{itemize}

Estas comprobaciones se ejecutaron con usuarios de prueba distintos, validando aislamiento por \texttt{auth.uid()} según lo especificado en el diseño de base de datos.

\subsection{Pruebas de usabilidad con usuarios objetivo}

Se realizó una evaluación cualitativa con dos usuarios externos (perfil entrenador junior) y un usuario interno del equipo. Se definieron tareas guiadas y se registró tiempo de ejecución, errores y percepción de dificultad.

\begin{longtable}{@{}p{4.5cm}p{2.5cm}p{2cm}p{4cm}@{}}
\toprule
\textbf{Tarea} & \textbf{Tiempo medio} & \textbf{Errores} & \textbf{Conclusión} \\
\midrule
\endhead
Registrar nuevo usuario & 52 s & Bajo & Flujo claro; mejorar pista visual en fecha \\
Crear rutina con 3 ejercicios & 2 min 18 s & Medio & Correcto; útil reforzar feedback al guardar \\
Buscar cliente y abrir detalle & 34 s & Bajo & Flujo intuitivo y directo \\
Enviar mensaje a cliente & 29 s & Bajo & Interacción comprensible y rápida \\
Consultar sesiones en calendario & 41 s & Bajo & Lectura temporal correcta \\
\bottomrule
\caption{Resumen de tareas de usabilidad y tiempos observados.}
\end{longtable}

La valoración global de uso fue positiva. Como mejoras propuestas, los participantes recomendaron incluir más microcopys de ayuda contextual en formularios y reforzar mensajes de confirmación tras acciones críticas.

\section{Trazabilidad entre requisitos y pruebas}

Para asegurar cobertura total, se relacionaron requisitos funcionales con los casos de prueba ejecutados:

\begin{longtable}{@{}p{2cm}p{8cm}p{4cm}@{}}
\toprule
\textbf{RF} & \textbf{Descripción resumida} & \textbf{Casos asociados} \\
\midrule
\endhead
RF-01 a RF-03 & Registro, login y roles & PT-01, PT-02, PT-03, PT-04, PT-05 \\
RF-04 a RF-05 & Resumen diario y próximas sesiones & PT-06, PNF-04 \\
RF-06 a RF-08 & Gestión y detalle de clientes & PT-06, PT-07, PT-08, PT-09 \\
RF-09 a RF-11 & Creación de rutinas y configuración de ejercicios & PT-10, PT-11, PT-12, PT-13 \\
RF-12 a RF-13 & Visualización en calendario & Tarea de usabilidad: calendario \\
RF-14 a RF-16 & Mensajería en tiempo real y no leídos & PT-14, PT-15, PT-16, PT-17, PNF-03 \\
\bottomrule
\caption{Matriz de trazabilidad Requisitos -- Pruebas.}
\end{longtable}

\section{Conclusión técnica de la validación}

El conjunto de evidencias funcionales y no funcionales permite concluir que FitLabs alcanza el nivel de calidad exigible para una entrega final de 2º DAM: funcionalidades principales completas, comportamiento estable en red real, seguridad de acceso correctamente aplicada y una usabilidad adecuada para el perfil objetivo.
